\documentclass[11pt]{article}
\usepackage[margin=1in]{geometry}
\usepackage{amsmath}
\usepackage{enumitem}
\title{Physics-Oriented Summary (Energy, Thermodynamics, and Kinetics)}
\date{2026-03-27}
\begin{document}
\maketitle

\section*{Why a Physics Summary Applies}
Although the uploaded set is chemistry-focused, several problems are built on physics principles: energy conservation, thermal response, exponential decay, and barrier-crossing dynamics.

\section*{Physics Concepts Extracted}
\begin{itemize}[leftmargin=*]
  \item \textbf{Energy landscapes}: reaction coordinate diagrams map potential energy vs. progress coordinate, analogous to barrier crossing in classical/statistical physics.
  \item \textbf{Activation barriers}: forward/reverse activation energies correspond to required energy for state transitions.
  \item \textbf{First-order dynamics}: exponential decay and logarithmic linearization are examples of dynamical systems with constant fractional change rate.
  \item \textbf{Thermal activation}: Arrhenius dependence links microscopic energy barrier $E_a$ to temperature.
  \item \textbf{Phase-change energetics}: heating curves illustrate latent heat, heat capacity regions, and constant-temperature transitions.
  \item \textbf{Entropy and microstates}: spontaneity trends rely on entropy changes tied to accessible microstate counts.
  \item \textbf{Gas partial pressure effects}: Henry-law usage reflects equilibrium between gas phase and dissolved phase under pressure constraints.
\end{itemize}

\section*{Key Equations in Physics Framing}
\[
\Delta G = \Delta H - T\Delta S
\]
\[
\ln k = -\frac{E_a}{R}\frac{1}{T} + \ln A
\]
\[
\ln [A]_t = -kt + \ln [A]_0
\]

\section*{Transferable Problem-Solving Heuristics}
\begin{enumerate}[leftmargin=*]
  \item Start with a clear sign convention for energy and rate quantities.
  \item Convert all values to coherent SI-like units before substitution.
  \item Use graph slope meanings explicitly ($-E_a/R$ for Arrhenius plots).
  \item For multi-species rates, enforce stoichiometric scaling factors first.
\end{enumerate}

\section*{Suggested Next Uploads for Deeper Physics Indexing}
\begin{itemize}[leftmargin=*]
  \item Dedicated mechanics/electricity worksheets.
  \item Wave/optics problems with derivations.
  \item Thermodynamics derivation notes (state functions, cycles, efficiency).
\end{itemize}

\end{document}
